% Part: methods
% Chapter: proofs
% Section: resources

\documentclass[../../../include/open-logic-section]{subfiles}

\begin{document}

\olfileid[hi]{mth}{prf}{res}

\section{अन्य संसाधन}

गणित में प्रमाण कैसे करें, इस विषय पर अनेक उपयोगी पुस्तकें हैं। विशेषतः
\emph{How to Read and do Proofs: An Introduction to Mathematical Thought
Processes} \citep{Solow2013} और \emph{How to Prove It: A Structured
Approach} \citep{Velleman2019} देखें।
\href{http://www.people.vcu.edu/~rhammack/BookOfProof/BookOfProof.pdf}{\emph{Book
of Proof}} \citep{Hammack2013} और
\href{https://scholarworks.gvsu.edu/books/7/}{\emph{Mathematical
Reasoning}} \citep{Sandstrum2019} प्रमाण पर ऐसी पुस्तकें हैं जो ऑनलाइन
निःशुल्क उपलब्ध हैं। दार्शनिकों को
\emph{More
Precisely: The Math you need to do Philosophy} \citep{Steinhart2018}
गणितीय तर्कणा की अच्छी प्रारंभिक पुस्तक लग सकती है।

इंटरनेट पर प्रमाणों की कई छोटी मार्गदर्शिकाएँ भी उपलब्ध हैं; उदाहरणतः,
\href{https://math.berkeley.edu/~hutching/teach/proofs.pdf}{``Introduction
  to Mathematical Arguments''} \citep{Hutchings2003} और
\href{https://eugeniacheng.com/wp-content/uploads/2017/02/cheng-proofguide.pdf}{``How to
  write proofs''} \citep{Cheng2004}.

\subsection{प्रेरक वीडियो}

क्या गृहकार्य करने की कोई प्रेरणा नहीं लग रही? मन उदास है? ये वीडियो सहायक
हो सकते हैं!

\begin{itemize}
\item \url{https://www.youtube.com/watch?v=ZXsQAXx_ao0}
\item \url{https://www.youtube.com/watch?v=BQ4yd2W50No}
\item \url{https://www.youtube.com/watch?v=StTqXEQ2l-Y}
\end{itemize}

\end{document}
